\cotefile{157-5}
\scanpage{325}\mspage{241}
% l. 1–5
Soit donc $\unsure{\xi|\zeta}_{\alpha}=(i_{\alpha},u'_{\alpha})$ \inserted{$(\alpha\in A,\operatorname{Card} A\leq\unsure{\pi|\kappa})$} une famille d'éléments de $i'\backslash I$;
\diagram{\begin{tikzcd}[ampersand replacement=\&,column sep=small]
 \& i'_{\alpha}\arrow[rr,bend left=12,"w_{\alpha}"]\arrow[d,no head] \& \deleted{\unsure{f}}j\arrow[r,"\deleted{\unsure{w}}","v"'] \& k\\
 i'\arrow[r,"u'_{\alpha}"]\arrow[rr,bend right=17,"v'"',no head] \& f(i\unsure{'}_{\alpha})\arrow[r,"u''_{\alpha}",no head]\& f(\unsure{j|j'})\arrow[r,"f(v)"]\& f(k)
\end{tikzcd}}
comme $I'$ satisfait à $\mathrm{PF}1_{\pi}$ (cf \unsure{4.9.6.9.7|4.9.6.9.2|4.9.6.9.1}), il existe $j'$ dans $I'$ et des flèches $u''_{\alpha}:f(i'_{\alpha})\to\deleted{\unsure{f(}}j'\deleted{\unsure{)}}$, telles que $u''_{\alpha}u'_{\alpha}=v'$ indépendamment de $\alpha$. \deleted{\unsure{v}} Comme $f$ satisfait
% l. 6–9
$\mathrm{F}1$, OPS que $j'$ est de la forme $f(j)$. \unsure{Appliquer|Appliquons} $\mathrm{F}3^{\bullet}_{\pi}$ aux \deleted{$u''_{\alpha}$ et} $i_{\alpha},j,u''_{\alpha}$, on \unsure{trouve|tr.} un $v:j\longrightarrow k$ et des $w_{\alpha}:\unsure{d_{\alpha}|i_{\alpha}}\unsure{\mathbin{-}|\longrightarrow}k$, tels que
\[
 \underbrace{f(v)u''_{\alpha}}_{\unsure{\text{indép.}|\text{nat}|\text{nât}}}=f(w_{\alpha})\qquad\forall\alpha\in A,
\]
% l. 10–13
d'où
\[
 \deleted{f(v)u''_{\alpha}}\quad f(w_{\alpha})u'_{\alpha}
 = f(v)\underbrace{u''_{\alpha}u'_{\alpha}}_{v'}=f(v)v'\quad\text{indép. des }\alpha,
\]
soit $\unsure{w'}$, donc les $w_{\alpha}$ définissent une flèche des $\xi_{\alpha}$ dans $\xi=(k,w')$, qed.

% l. 14–18
$b''_{\pi}\Longrightarrow\unsure{b}_{\pi}$ \inserted{si $f$ est fidèle}. On est \unsure{ramené} à prouver que si $b''_{\pi}$ est satisfait, \inserted{$f$ fidèle,} alors \deleted{\unsure{on a la condition} $\mathrm{F}3$.} \inserted{les} catégories $i'\backslash I$, dont on sait déjà qu'elles satisfont $\mathrm{PF}1_{\pi}$, satisfont également $\mathrm{PF}2_{\pi}$.

% l. 19–26
\diagram{\begin{tikzcd}[ampersand replacement=\&,column sep=small]
 \& i\deleted{\unsure{\alpha}}\arrow[r,shift left=2,"u_{\alpha}"]\arrow[r]\arrow[r,shift right=2]\arrow[d,no head]\& j\arrow[r,"\unsure{v.}",dotted]\& k\\
 i'\arrow[r,"u'"]\arrow[rr,bend right=15,"v'"',no head]\& f(i\unsure{_{\alpha}|{}})\arrow[r,shift left=2]\arrow[r]\arrow[r,shift right=2]\& f(j)\arrow[r]\& f(k)
\end{tikzcd}}
Soient donc $\xi\unsure{_{1}}=(i,u')$, $\eta=(j,v')$ deux objets de $i'\backslash I$, et $(u_{\alpha})_{\alpha\in A}$ une famille de flèches $\xi\longrightarrow\eta$ \inserted{(avec $\operatorname{Card} A\leq\pi$)}, donc des flèches dans $I$ telles que
\[
 \underline{f(u_{\alpha})u'=v'}\qquad\forall\alpha\unsure{\in}A.
\]
Comme $I'$ est \unsure{grde|grdes|gr des|qv des} devant $\pi$, on peut égaliser par $f(j)\longrightarrow k'$, et comme \deleted{$\unsure{f}$} $k'\backslash I\ne\varnothing$, OPS $k'$ de la forme $f(k)$. Or
% archive sheet number: 324
